\section{Higher codimension gysin homomorphisms}
\label{section-gysin-higher-codimension}

\noindent
Let $(S, \delta)$ be as in Situation \ref{situation-setup}. Let $X$ be a scheme
locally of finite type over $S$. In this section we are going to consider
triples
$$
(Z \to X, \mathcal{N}, \sigma : \mathcal{N}^\vee \to \mathcal{C}_{Z/X})
$$
consisting of a closed immersion $Z \to X$ and a locally free
$\mathcal{O}_Z$-module $\mathcal{N}$ and a surjection
$\sigma : \mathcal{N}^\vee \to \mathcal{C}_{Z/X}$ from the dual
of $\mathcal{N}$ to the conormal sheaf of $Z$ in $X$, see
Morphisms, Section \ref{morphisms-section-conormal-sheaf}.
We will say
$\mathcal{N}$ is a {\it virtual normal sheaf for $Z$ in $X$}.

\begin{lemma}
\label{lemma-pullback-virtual-normal-sheaf}
Let $(S, \delta)$ be as in Situation \ref{situation-setup}. Let
$$
\xymatrix{
Z' \ar[r] \ar[d]_g & X' \ar[d]^f \\
Z \ar[r] & X
}
$$
be a cartesian diagram of schemes locally of finite type over $S$
whose horizontal arrows are closed immersions.
If $\mathcal{N}$ is a virtual normal sheaf for $Z$ in $X$, then
$\mathcal{N}' = g^*\mathcal{N}$ is a virtual normal sheaf for
$Z'$ in $X'$.
\end{lemma}

\begin{proof}
This follows from the surjectivity of the map
$g^*\mathcal{C}_{Z/X} \to \mathcal{C}_{Z'/X'}$ proved in
Morphisms, Lemma \ref{morphisms-lemma-conormal-functorial-flat}.
\end{proof}

\noindent
Let $(S, \delta)$ be as in Situation \ref{situation-setup}. Let $X$ be a scheme
locally of finite type over $S$. Let $\mathcal{N}$ be a virtual normal bundle
for a closed immersion $Z \to X$. In this situation we set
$$
p : N = \underline{\Spec}_Z(\text{Sym}(\mathcal{N}^\vee)) \longrightarrow Z
$$
equal to the vector bundle over $Z$
whose sections correspond to sections of $\mathcal{N}$.
In this situation we have canonical closed immersions
$$
C_ZX \longrightarrow N_ZX \longrightarrow N
$$
The first closed immersion is Divisors, Equation
(\ref{divisors-equation-normal-cone-in-normal-bundle})
and the second closed immersion corresponds to the surjection
$\text{Sym}(\mathcal{N}^\vee) \to \text{Sym}(\mathcal{C}_{Z/X})$
induced by $\sigma$.
Let
$$
b : W \longrightarrow \mathbf{P}^1_X
$$
be the blowing up in $\infty(Z)$ constructed in
Section \ref{section-blowup-Z-first}. By
Lemma \ref{lemma-gysin-at-infty}
we have a canonical bivariant class in
$$
C \in A^0(W_\infty \to X)
$$
Consider the open immersion $j : C_ZX \to W_\infty$ of
(\ref{item-cone-is-open}) and the closed immersion
$i : C_ZX \to N$ constructed above. By Lemma \ref{lemma-vectorbundle}
for every $\alpha \in \CH_k(X)$ there exists a unique
$\beta \in \CH_*(Z)$ such that
$$
i_*j^*(C \cap \alpha) = p^*\beta
$$
We set $c(Z \to X, \mathcal{N}) \cap \alpha = \beta$.

\begin{lemma}
\label{lemma-construction-gysin}
The construction above defines a bivariant class\footnote{The
notation $A^*(Z \to X)^\wedge$ is discussed in
Remark \ref{remark-completion-bivariant}.
If $X$ is quasi-compact, then $A^*(Z \to X)^\wedge = A^*(Z \to X)$.}
$$
c(Z \to X, \mathcal{N}) \in A^*(Z \to X)^\wedge
$$
and moreover the construction is compatible with base change
as in Lemma \ref{lemma-pullback-virtual-normal-sheaf}.
If $\mathcal{N}$ has constant rank $r$, then
$c(Z \to X, \mathcal{N}) \in A^r(Z \to X)$.
\end{lemma}

\begin{proof}
Since both $i_* \circ j^* \circ C$ and $p^*$ are bivariant classes
(see Lemmas \ref{lemma-flat-pullback-bivariant} and
\ref{lemma-push-proper-bivariant}) we can use the equation
$$
i_* \circ j^* \circ C  = p^* \circ c(Z \to X, \mathcal{N})
$$
(suitably interpreted) to define $c(Z \to X, \mathcal{N})$
as a bivariant class. This works because $p^*$ is always
bijective on chow groups by Lemma \ref{lemma-vectorbundle}.

\medskip\noindent
Let $X' \to X$, $Z' \to X'$, and $\mathcal{N}'$ be as in
Lemma \ref{lemma-pullback-virtual-normal-sheaf}. Write
$c = c(Z \to X, \mathcal{N})$ and $c' = c(Z' \to X', \mathcal{N}')$.
The second statement of the lemma means that $c'$ is the restriction of $c$
as in Remark \ref{remark-restriction-bivariant}. Since we claim this
is true for all $X'/X$ locally of finite type, a formal argument
shows that it suffices to check that $c' \cap \alpha' = c \cap \alpha'$
for $\alpha' \in \CH_k(X')$.
To see this, note that we have a commutative diagram
$$
\xymatrix{
C_{Z'}X' \ar[d] \ar[r] &
W'_\infty \ar[d] \ar[r] &
W' \ar[d] \ar[r] &
\mathbf{P}^1_{X'} \ar[d] \\
C_ZX \ar[r] &
W_\infty \ar[r] &
W \ar[r] &
\mathbf{P}^1_X
}
$$
which induces closed immersions:
$$
W' \to W \times_{\mathbf{P}^1_X} \mathbf{P}^1_{X'},\quad
W'_\infty \to W_\infty \times_X X',\quad
C_{Z'}X' \to C_ZX \times_Z Z'
$$
To get $c \cap \alpha'$ we use the class $C \cap \alpha'$
defined using the morphism
$W \times_{\mathbf{P}^1_X} \mathbf{P}^1_{X'} \to \mathbf{P}^1_{X'}$
in Lemma \ref{lemma-gysin-at-infty}.
To get $c' \cap \alpha'$ on the other hand, we use the class
$C' \cap \alpha'$ defined using the morphism $W' \to \mathbf{P}^1_{X'}$.
By Lemma \ref{lemma-gysin-at-infty-independent} the pushforward of
$C' \cap \alpha'$ by the closed immersion
$W'_\infty \to (W \times_{\mathbf{P}^1_X} \mathbf{P}^1_{X'})_\infty$,
is equal to $C \cap \alpha'$. Hence the same is true for the pullbacks
to the opens
$$
C_{Z'}X' \subset W'_\infty,\quad
C_ZX \times_Z Z' \subset (W \times_{\mathbf{P}^1_X} \mathbf{P}^1_{X'})_\infty
$$
by Lemma \ref{lemma-flat-pullback-proper-pushforward}.
Since we have a commutative diagram
$$
\xymatrix{
C_{Z'} X' \ar[d] \ar[r] & N' \ar@{=}[d] \\
C_ZX \times_Z Z' \ar[r] & N \times_Z Z'
}
$$
these classes pushforward to the same class on $N'$ which
proves that we obtain the same element $c \cap \alpha' = c' \cap \alpha'$
in $\CH_*(Z')$.
\end{proof}

\begin{lemma}
\label{lemma-gysin-decompose}
Let $(S, \delta)$ be as in Situation \ref{situation-setup}. Let $X$ be a scheme
locally of finite type over $S$. Let $\mathcal{N}$ be a virtual normal
sheaf for a closed subscheme $Z$ of $X$. Suppose that we have a short
exact sequence $0 \to \mathcal{N}' \to \mathcal{N} \to \mathcal{E} \to 0$
of finite locally free $\mathcal{O}_Z$-modules such that the given surjection
$\sigma : \mathcal{N}^\vee \to \mathcal{C}_{Z/X}$ factors through a map
$\sigma' : (\mathcal{N}')^\vee \to \mathcal{C}_{Z/X}$.
Then
$$
c(Z \to X, \mathcal{N}) = c_{top}(\mathcal{E}) \circ c(Z \to X, \mathcal{N}')
$$
as bivariant classes.
\end{lemma}

\begin{proof}
Denote $N' \to N$ the closed immersion of vector bundles corresponding
to the surjection $\mathcal{N}^\vee \to (\mathcal{N}')^\vee$. Then we
have closed immersions
$$
C_ZX \to N' \to N
$$
Thus the desired relationship between the bivariant classes follows
immediately from Lemma \ref{lemma-easy-virtual-class}.
\end{proof}

\begin{lemma}
\label{lemma-gysin-excess}
Let $(S, \delta)$ be as in Situation \ref{situation-setup}. Consider
a cartesian diagram
$$
\xymatrix{
Z' \ar[r] \ar[d]_g & X' \ar[d]^f \\
Z \ar[r] & X
}
$$
of schemes locally of finite type over $S$ whose horizontal arrows
are closed immersions. Let $\mathcal{N}$, resp.\ $\mathcal{N}'$
be a virtual normal sheaf for $Z \subset X$, resp.\ $Z' \to X'$.
Assume given a short exact sequence
$0 \to \mathcal{N}' \to g^*\mathcal{N} \to \mathcal{E} \to 0$
of finite locally free modules on $Z'$ such that the diagram
$$
\xymatrix{
g^*\mathcal{N}^\vee \ar[r] \ar[d] &
(\mathcal{N}')^\vee \ar[d] \\
g^*\mathcal{C}_{Z/X} \ar[r] &
\mathcal{C}_{Z'/X'}
}
$$
commutes. Then we have
$$
res(c(Z \to X, \mathcal{N})) =
c_{top}(\mathcal{E}) \circ c(Z' \to X', \mathcal{N}')
$$
in $A^*(Z' \to X')^\wedge$.
\end{lemma}

\begin{proof}
By Lemma \ref{lemma-construction-gysin} we have
$res(c(Z \to X, \mathcal{N})) = c(Z' \to X', g^*\mathcal{N})$
and the equality follows from Lemma \ref{lemma-gysin-decompose}.
\end{proof}

\noindent
Let $(S, \delta)$ be as in Situation \ref{situation-setup}. Let $X$ be a scheme
locally of finite type over $S$. Let $\mathcal{N}$ be a virtual normal
sheaf for a closed subscheme $Z$ of $X$. Let $Y \to X$ be a morphism
which is locally of finite type. Assume $Z \times_X Y \to Y$ is a
regular closed immersion, see
Divisors, Section \ref{divisors-section-regular-immersions}.
In this case the conormal sheaf $\mathcal{C}_{Z \times_X Y/Y}$ is a finite
locally free $\mathcal{O}_{Z \times_X Y}$-module and we obtain a short
exact sequence
$$
0 \to \mathcal{E}^\vee \to
\mathcal{N}^\vee|_{Z \times_X Y} \to \mathcal{C}_{Z \times_X Y/Y} \to 0
$$
The quotient $\mathcal{N}|_{Y \times_X Z} \to \mathcal{E}$ is called the
{\it excess normal sheaf} of the situation.

\begin{lemma}
\label{lemma-gysin-fundamental}
In the situation described just above assume $\dim_\delta(Y) = n$
and that $\mathcal{C}_{Y \times_X Z/Z}$ has constant rank $r$.
Then
$$
c(Z \to X, \mathcal{N}) \cap [Y]_n =
c_{top}(\mathcal{E}) \cap [Z \times_X Y]_{n - r}
$$
in $\CH_*(Z \times_X Y)$.
\end{lemma}

\begin{proof}
The bivariant class $c_{top}(\mathcal{E}) \in A^*(Z \times_X Y)$ was
defined in Remark \ref{remark-top-chern-class}.
By Lemma \ref{lemma-construction-gysin} we may replace $X$ by $Y$.
Thus we may assume $Z \to X$ is a regular closed immersion
of codimension $r$, we have $\dim_\delta(X) = n$, and we have
to show that $c(Z \to X, \mathcal{N}) \cap [X]_n =
c_{top}(\mathcal{E}) \cap [Z]_{n - r}$ in $\CH_*(Z)$.
By Lemma \ref{lemma-gysin-decompose} we may even assume
$\mathcal{N}^\vee \to \mathcal{C}_{Z/X}$ is an isomorphism.
In other words, we have to show
$c(Z \to X, \mathcal{C}_{Z/X}^\vee) \cap [X]_n = [Z]_{n - r}$ in $\CH_*(Z)$.

\medskip\noindent
Let us trace through the steps in the definition of
$c(Z \to X, \mathcal{C}_{Z/X}^\vee) \cap [X]_n$. Let
$b : W \to \mathbf{P}^1_X$
be the blowing up of $\infty(Z)$. We first have to compute
$C \cap [X]_n$ where $C \in A^0(W_\infty \to X)$ is
the class of Lemma \ref{lemma-gysin-at-infty}.
To do this, note that $[W]_{n + 1}$
is a cycle on $W$ whose restriction to $\mathbf{A}^1_X$ is
equal to the flat pullback of $[X]_n$. Hence $C \cap [X]_n$
is equal to $i_\infty^*[W]_{n + 1}$. Since $W_\infty$ is an
effective Cartier divisor on $W$ we have
$i_\infty^*[W]_{n + 1} = [W_\infty]_n$, see Lemma \ref{lemma-easy-gysin}.
The restriction of this class to the open $C_ZX \subset W_\infty$
is of course just $[C_ZX]_n$. Because $Z \subset X$ is regularly
embedded we have
$$
\mathcal{C}_{Z/X, *} = \text{Sym}(\mathcal{C}_{Z/X})
$$
as graded $\mathcal{O}_Z$-algebras, see
Divisors, Lemma \ref{divisors-lemma-quasi-regular-immersion}.
Hence $p : N = C_ZX \to Z$ is the structure morphism of the
vector bundle associated to the finite locally free module
$\mathcal{C}_{Z/X}$ of rank $r$. Then it is clear that
$p^*[Z]_{n - r} = [C_ZX]_n$ and the proof is complete.
\end{proof}

\begin{lemma}
\label{lemma-gysin-easy}
Let $(S, \delta)$ be as in Situation \ref{situation-setup}. Let $X$ be a scheme
locally of finite type over $S$. Let $\mathcal{N}$ be a virtual normal
sheaf for a closed subscheme $Z$ of $X$. Let $Y \to X$ be a morphism
which is locally of finite type. Given integers $r$, $n$ assume
\begin{enumerate}
\item $\mathcal{N}$ is locally free of rank $r$,
\item every irreducible component of $Y$ has $\delta$-dimension $n$,
\item $\dim_\delta(Z \times_X Y) \leq n - r$, and
\item for $\xi \in Z \times_X Y$ with $\delta(\xi) = n - r$
the local ring $\mathcal{O}_{Y, \xi}$ is Cohen-Macaulay.
\end{enumerate}
Then $c(Z \to X, \mathcal{N}) \cap [Y]_n = [Z \times_X Y]_{n - r}$
in $\CH_{n - r}(Z \times_X Y)$.
\end{lemma}

\begin{proof}
The statement makes sense as $Z \times_X Y$ is a closed subscheme of $Y$.
Because $\mathcal{N}$ has rank $r$ we know that
$c(Z \to X, \mathcal{N}) \cap [Y]_n$ is in $\CH_{n - r}(Z \times_X Y)$.
Since $\dim_\delta(Z \cap Y) \leq n - r$ the chow group
$\CH_{n - r}(Z \times_X Y)$ is freely generated by the
cycle classes of the irreducible components $W \subset Z \times_X Y$
of $\delta$-dimension $n - r$. Let $\xi \in W$ be the generic point.
By assumption (2) we see that $\dim(\mathcal{O}_{Y, \xi}) = r$.
On the other hand, since $\mathcal{N}$ has rank $r$ and since
$\mathcal{N}^\vee \to \mathcal{C}_{Z/X}$ is surjective, we see that
the ideal sheaf of $Z$ is locally cut out by $r$ equations.
Hence the quasi-coherent ideal sheaf $\mathcal{I} \subset \mathcal{O}_Y$
of $Z \times_X Y$ in $Y$ is locally generated by $r$ elements.
Since $\mathcal{O}_{Y, \xi}$ is Cohen-Macaulay of dimension $r$
and since $\mathcal{I}_\xi$ is an ideal of definition (as $\xi$ is
a generic point of $Z \times_X Y$) it follows that $\mathcal{I}_\xi$
is generated by a regular sequence
(Algebra, Lemma \ref{algebra-lemma-reformulate-CM}).
By Divisors, Lemma \ref{divisors-lemma-Noetherian-scheme-regular-ideal}
we see that $\mathcal{I}$ is generated by a regular sequence over
an open neighbourhood $V \subset Y$ of $\xi$. By our description of
$\CH_{n - r}(Z \times_X Y)$ it suffices to show that
$c(Z \to X, \mathcal{N}) \cap [V]_n = [Z \times_X V]_{n - r}$
in $\CH_{n - r}(Z \times_X V)$. This follows from
Lemma \ref{lemma-gysin-fundamental}
because the excess normal sheaf is $0$ over $V$.
\end{proof}

\begin{lemma}
\label{lemma-gysin-agrees}
Let $(S, \delta)$ be as in Situation \ref{situation-setup}. Let $X$ be a scheme
locally of finite type over $S$. Let $(\mathcal{L}, s, i : D \to X)$
be a triple as in Definition \ref{definition-gysin-homomorphism}.
The Gysin homomorphism $i^*$ viewed as an element of $A^1(D \to X)$
(see Lemma \ref{lemma-gysin-bivariant}) is the same as the bivariant class
$c(D \to X, \mathcal{N}) \in A^1(D \to X)$
constructed using $\mathcal{N} = i^*\mathcal{L}$
viewed as a virtual normal sheaf for $D$ in $X$.
\end{lemma}

\begin{proof}
We will use the criterion of Lemma \ref{lemma-bivariant-zero}.
Thus we may assume that $X$ is an integral scheme and
we have to show that $i^*[X]$ is equal to $c \cap [X]$.
Let $n = \dim_\delta(X)$. As usual, there are two cases.

\medskip\noindent
If $X = D$, then we see that both classes are represented by
$c_1(\mathcal{N}) \cap [X]_n$. See Lemma \ref{lemma-gysin-fundamental}
and Definition \ref{definition-gysin-homomorphism}.

\medskip\noindent
If $D \not = X$, then $D \to X$ is an effective Cartier divisor
and in particular a regular closed immersion of codimension $1$.
Again by Lemma \ref{lemma-gysin-fundamental} we conclude
$c(D \to X, \mathcal{N}) \cap [X]_n = [D]_{n - 1}$. The same
is true by definition for the Gysin homomorphism and we conclude
once again.
\end{proof}

\begin{lemma}
\label{lemma-gysin-commutes}
Let $(S, \delta)$ be as in Situation \ref{situation-setup}. Let $X$ be a scheme
locally of finite type over $S$. Let $Z \subset X$ be a closed subscheme
with virtual normal sheaf $\mathcal{N}$. Let $Y \to X$ be locally of
finite type and $c \in A^*(Y \to X)$. Then $c$ and $c(Z \to X, \mathcal{N})$
commute (Remark \ref{remark-bivariant-commute}).
\end{lemma}

\begin{proof}
To check this we may use Lemma \ref{lemma-bivariant-zero}.
Thus we may assume $X$ is an integral scheme and we have to show
$c \cap c(Z \to X, \mathcal{N}) \cap [X] =
c(Z \to X, \mathcal{N}) \cap c \cap [X]$ in $\CH_*(Z \times_X Y)$.

\medskip\noindent
If $Z = X$, then $c(Z \to X, \mathcal{N}) = c_{top}(\mathcal{N})$ by
Lemma \ref{lemma-gysin-fundamental} which commutes
with the bivariant class $c$, see Lemma \ref{lemma-cap-commutative-chern}.

\medskip\noindent
Assume that $Z$ is not equal to $X$. By Lemma \ref{lemma-bivariant-zero}
it even suffices to prove the result after blowing up $X$ (in a nonzero ideal).
Let us blowup $X$ in the ideal sheaf of $Z$. This reduces us to the case
where $Z$ is an effective Cartier divisor, see
Divisors, Lemma
\ref{divisors-lemma-blowing-up-gives-effective-Cartier-divisor},

\medskip\noindent
If $Z$ is an effective Cartier divisor, then we have
$$
c(Z \to X, \mathcal{N}) =
c_{top}(\mathcal{E}) \circ i^*
$$
where $i^* \in A^1(Z \to X)$ is the Gysin homomorphism
associated to $i : Z \to X$ (Lemma \ref{lemma-gysin-bivariant})
and $\mathcal{E}$ is the dual of the kernel of
$\mathcal{N}^\vee \to \mathcal{C}_{Z/X}$, see
Lemmas \ref{lemma-gysin-decompose} and \ref{lemma-gysin-agrees}.
Then we conclude because Chern classes are in the center of the
bivariant ring (in the strong sense formulated in
Lemma \ref{lemma-cap-commutative-chern}) and $c$ commutes
with the Gysin homomorphism $i^*$ by definition of bivariant classes.
\end{proof}

\noindent
Let $(S, \delta)$ be as in Situation \ref{situation-setup}. Let $X$ be an
integral scheme locally of finite type over $S$ of $\delta$-dimension $n$.
Let $Z \subset Y \subset X$ be closed subschemes which are both effective
Cartier divisors in $X$. Denote $o : Y \to C_Y X$ the zero section of the
normal line cone of $Y$ in $X$. As $C_YX$ is a line bundle over $Y$
we obtain a bivariant class $o^* \in A^1(Y \to C_YX)$, see
Lemma \ref{lemma-gysin-bivariant}.

\begin{lemma}
\label{lemma-relation-normal-cones}
With notation as above we have
$$
o^*[C_ZX]_n = [C_Z Y]_{n - 1}
$$
in $\CH_{n - 1}(Y \times_{o, C_Y X} C_ZX)$.
\end{lemma}

\begin{proof}
Denote $W \to \mathbf{P}^1_X$ the blowing up of $\infty(Z)$ as in
Section \ref{section-blowup-Z-first}.
Similarly, denote $W' \to \mathbf{P}^1_X$ the blowing up of $\infty(Y)$.
Since $\infty(Z) \subset \infty(Y)$ we get an opposite inclusion
of ideal sheaves and hence a map of the graded algebras
defining these blowups. This produces a rational morphism from $W$
to $W'$ which in fact has a canonical representative
$$
W \supset U \longrightarrow W'
$$
See Constructions, Lemma \ref{constructions-lemma-morphism-relative-proj}.
A local calculation (omitted) shows that $U$ contains at least all points
of $W$ not lying over $\infty$ and the open subscheme $C_Z X$ of the special
fibre. After shrinking $U$ we may assume $U_\infty = C_Z X$ and
$\mathbf{A}^1_X \subset U$. Another local calculation (omitted)
shows that the morphism $U_\infty \to W'_\infty$
induces the canonical morphism $C_Z X \to C_Y X \subset W'_\infty$
of normal cones induced by the inclusion of ideals sheaves
coming from $Z \subset Y$. Denote $W'' \subset W$ the strict transform of
$\mathbf{P}^1_Y \subset \mathbf{P}^1_X$ in $W$. Then $W''$ is the blowing
up of $\mathbf{P}^1_Y$ in $\infty(Z)$ by
Divisors, Lemma \ref{divisors-lemma-strict-transform}
and hence $(W'' \cap U)_\infty = C_ZY$.

\medskip\noindent
Consider the effective Cartier divisor $i : \mathbf{P}^1_Y \to W'$
from (\ref{item-find-Z-in-blowup}) and its associated bivariant class
$i^* \in A^1(\mathbf{P}^1_Y \to W')$ from Lemma \ref{lemma-gysin-bivariant}.
We similarly denote $(i'_\infty)^* \in A^1(W'_\infty \to W')$ the
Gysin map at infinity. Observe that the restriction of $i'_\infty$
(Remark \ref{remark-restriction-bivariant}) to $U$ is the restriction of
$i_\infty^* \in A^1(W_\infty \to W)$ to $U$. On the one hand we have
$$
(i'_\infty)^* i^* [U]_{n + 1} =
i_\infty^* i^* [U]_{n + 1} =
i_\infty^* [(W'' \cap U)_\infty]_{n + 1} =
[C_ZY]_n
$$
because $i_\infty^*$ kills all classes supported over $\infty$, because
$i^*[U]$ and $[W'']$ agree as cycles over $\mathbf{A}^1$, and because
$C_ZY$ is the fibre of $W'' \cap U$ over $\infty$.
On the other hand, we have
$$
(i'_\infty)^* i^* [U]_{n + 1} =
i^* i_\infty^*[U]_{n + 1} =
i^* [U_\infty] =
o^*[C_YX]_n
$$
because $(i'_\infty)^*$ and $i^*$ commute
(Lemma \ref{lemma-gysin-commutes-gysin})
and because the fibre of $i : \mathbf{P}^1_Y \to W'$ over $\infty$
factors as $o : Y \to C_YX$ and the open immersion $C_YX \to W'_\infty$.
The lemma follows.
\end{proof}

\begin{lemma}
\label{lemma-gysin-composition}
Let $(S, \delta)$ be as in Situation \ref{situation-setup}.
Let $Z \subset Y \subset X$ be closed subschemes of a scheme locally
of finite type over $S$.
Let $\mathcal{N}$ be a virtual normal sheaf for $Z \subset X$.
Let $\mathcal{N}'$ be a virtual normal sheaf for $Z \subset Y$.
Let $\mathcal{N}''$ be a virtual normal sheaf for $Y \subset X$.
Assume there is a commutative diagram
$$
\xymatrix{
(\mathcal{N}'')^\vee|_Z \ar[r] \ar[d] &
\mathcal{N}^\vee \ar[r] \ar[d] &
(\mathcal{N}')^\vee \ar[d] \\
\mathcal{C}_{Y/X}|_Z \ar[r] &
\mathcal{C}_{Z/X} \ar[r] &
\mathcal{C}_{Z/Y}
}
$$
where the sequence at the bottom is from More on Morphisms, Lemma
\ref{more-morphisms-lemma-transitivity-conormal} and the top
sequence is a short exact sequence. Then
$$
c(Z \to X, \mathcal{N}) =
c(Z \to Y, \mathcal{N}') \circ c(Y \to X, \mathcal{N}'')
$$
in $A^*(Z \to X)^\wedge$.
\end{lemma}

\begin{proof}
Observe that the assumptions remain satisfied after any base change
by a morphism $X' \to X$ which is locally of finite type (the short
exact sequence of virtual normal sheaves is locally split hence
remains exact after any base change). Thus to check the
equality of bivariant classes we may use Lemma \ref{lemma-bivariant-zero}.
Thus we may assume $X$ is an integral scheme and we have to show
$c(Z \to X, \mathcal{N}) \cap [X] =
c(Z \to Y, \mathcal{N}') \cap c(Y \to X, \mathcal{N}'') \cap [X]$.

\medskip\noindent
If $Y = X$, then we have
\begin{align*}
c(Z \to Y, \mathcal{N}') \cap c(Y \to X, \mathcal{N}'') \cap [X]
& =
c(Z \to Y, \mathcal{N}') \cap c_{top}(\mathcal{N}'') \cap [Y] \\
& =
c_{top}(\mathcal{N}''|_Z) \cap c(Z \to Y, \mathcal{N}') \cap [Y] \\
& =
c(Z \to X, \mathcal{N}) \cap [X] 
\end{align*}
The first equality by Lemma \ref{lemma-gysin-decompose}.
The second because Chern classes commute with bivariant classes
(Lemma \ref{lemma-cap-commutative-chern}).
The third equality by Lemma \ref{lemma-gysin-decompose}.

\medskip\noindent
Assume $Y \not = X$. By Lemma \ref{lemma-bivariant-zero}
it even suffices to prove the result after blowing up $X$ in a nonzero ideal.
Let us blowup $X$ in the product of the ideal sheaf of $Y$ and the ideal
sheaf of $Z$. This reduces us to the case where both $Y$ and $Z$ are
effective Cartier divisors on $X$, see
Divisors, Lemmas
\ref{divisors-lemma-blowing-up-gives-effective-Cartier-divisor} and
\ref{divisors-lemma-blowing-up-two-ideals}.

\medskip\noindent
Denote $\mathcal{N}'' \to \mathcal{E}$ the surjection of finite locally
free $\mathcal{O}_Z$-modules such that
$0 \to \mathcal{E}^\vee \to (\mathcal{N}'')^\vee \to \mathcal{C}_{Y/X} \to 0$
is a short exact sequence. Then $\mathcal{N} \to \mathcal{E}|_Z$
is a surjection as well. Denote $\mathcal{N}_1$ the finite locally free kernel
of this map and observe that $\mathcal{N}^\vee \to \mathcal{C}_{Z/X}$
factors through $\mathcal{N}_1$.
By Lemma \ref{lemma-gysin-decompose} we have
$$
c(Y \to X, \mathcal{N}'') = c_{top}(\mathcal{E}) \circ
c(Y \to X, \mathcal{C}_{Y/X}^\vee)
$$
and
$$
c(Z \to X, \mathcal{N}) = c_{top}(\mathcal{E}|_Z) \circ
c(Z \to X, \mathcal{N}_1)
$$
Since Chern classes of bundles commute with bivariant classes
(Lemma \ref{lemma-cap-commutative-chern})
it suffices to prove
$$
c(Z \to X, \mathcal{N}_1) =
c(Z \to Y, \mathcal{N}') \circ c(Y \to X, \mathcal{C}_{Y/X}^\vee)
$$
in $A^*(Z \to X)$. This we may assume that $\mathcal{N}'' = \mathcal{C}_{Y/X}$.
This reduces us to the case discussed in the next paragraph.

\medskip\noindent
In this paragraph $Z$ and $Y$ are effective Cartier divisors on $X$
integral of dimension $n$, we have $\mathcal{N}'' = \mathcal{C}_{Y/X}$.
In this case $c(Y \to X, \mathcal{C}_{Y/X}^\vee) \cap [X] = [Y]_{n - 1}$ by
Lemma \ref{lemma-gysin-fundamental}. Thus we have to prove that
$c(Z \to X, \mathcal{N}) \cap [X] = c(Z \to Y, \mathcal{N}') \cap [Y]_{n - 1}$.
Denote $N$ and $N'$ the vector bundles over $Z$ associated to
$\mathcal{N}$ and $\mathcal{N}'$. Consider the commutative diagram
$$
\xymatrix{
N' \ar[r]_i &
N \ar[r] &
(C_Y X) \times_Y Z \\
C_Z Y \ar[r] \ar[u] &
C_Z X \ar[u]
}
$$
of cones and vector bundles over $Z$. Observe that $N'$ is a relative
effective Cartier divisor in $N$ over $Z$ and that
$$
\xymatrix{
N' \ar[d] \ar[r]_i & N \ar[d] \\
Z \ar[r]^-o & (C_Y X) \times_Y Z
}
$$
is cartesian where $o$ is the zero section of the line bundle
$C_Y X$ over $Y$. By
Lemma \ref{lemma-relation-normal-cones} we have $o^*[C_ZX]_n = [C_Z Y]_{n - 1}$
in
$$
\CH_{n - 1}(Y \times_{o, C_Y X} C_ZX) =
\CH_{n - 1}(Z \times_{o, (C_Y X) \times_Y Z} C_ZX)
$$
By the cartesian property of
the square above this implies that
$$
i^*[C_ZX]_n = [C_Z Y]_{n - 1}
$$
in $\CH_{n - 1}(N')$. Now observe that
$\gamma = c(Z \to X, \mathcal{N}) \cap [X]$ and
$\gamma' = c(Z \to Y, \mathcal{N}') \cap [Y]_{n - 1}$
are characterized by $p^*\gamma = [C_Z X]_n$ in $\CH_n(N)$
and by $(p')^*\gamma' = [C_Z Y]_{n - 1}$ in $\CH_{n - 1}(N')$.
Hence the proof is finished as $i^* \circ p^* = (p')^*$ by
Lemma \ref{lemma-relative-effective-cartier}.
\end{proof}

\begin{remark}[Variant for immersions]
\label{remark-gysin-for-immersion}
Let $(S, \delta)$ be as in Situation \ref{situation-setup}.
Let $X$ be a scheme locally of finite type over $S$.
Let $i : Z \to X$ be an immersion of schemes.
In this situation
\begin{enumerate}
\item the conormal sheaf $\mathcal{C}_{Z/X}$
of $Z$ in $X$ is defined
(Morphisms, Definition \ref{morphisms-definition-conormal-sheaf}),
\item we say a pair consisting of a finite locally free $\mathcal{O}_Z$-module
$\mathcal{N}$ and a surjection $\sigma : \mathcal{N}^\vee \to \mathcal{C}_{Z/X}$
is a virtual normal bundle for the immersion $Z \to X$,
\item choose an open subscheme $U \subset X$ such that $Z \to X$
factors through a closed immersion $Z \to U$ and set
$c(Z \to X, \mathcal{N}) = c(Z \to U, \mathcal{N}) \circ (U \to X)^*$.
\end{enumerate}
The bivariant class $c(Z \to X, \mathcal{N})$ does not depend on the choice
of the open subscheme $U$. All of the lemmas have immediate counterparts
for this slightly more general construction. We omit the details.
\end{remark}







